\chapter{The Math Toolkit}

\section{Big and small numbers (scientific notation)}

\begin{kid}
Light is \emph{so} fast that it could go around the Earth about seven times in one second. To write such big numbers --- and tiny ones like ``one billionth of a second'' --- scientists use a shortcut: count the zeros and write that count as a little number up high. $10^3$ means 1 followed by 3 zeros = 1000. $10^{-3}$ means ``1 divided into 1000 pieces'' = 0.001.
\end{kid}

\section{Angles: degrees and radians}

\fig[0.6]{radians}{A radian is the angle whose arc is exactly one radius long.}

\begin{kid}
Imagine a pizza. If you cut a slice whose crust is exactly as long as the line from the center to the crust, that slice's angle is \textbf{one radian}. A whole pizza has a bit more than 6 of those slices (6.28, which is $2\pi$).
\end{kid}

\section{Sine and cosine}

\fig[0.38]{unitcircle}{On a circle of radius 1, cosine is the ``across'' part and sine is the ``up'' part.}

\begin{kid}
Lean a 1-meter stick against a wall. The bottom of the stick is some distance from the wall (that is \textbf{cosine}) and the top is some height up the wall (that is \textbf{sine}). If the stick leans a lot, cosine is big and sine is small. If it stands almost straight up, sine is big and cosine is small.
\end{kid}

\section{Vectors}

\begin{kid}
A vector is an arrow: it says \emph{how far} and \emph{which way}. ``Walk 3 steps east and 1 step north'' is a vector. Two walks one after the other add up to one bigger walk.
\end{kid}

\section{Statistics: mean, standard deviation, RMSE}

\fig[0.75]{gaussian}{Random measurement errors usually follow a bell curve. About 68\% fall within one standard deviation $\sigma$ and 95\% within 1.96$\sigma$.}

\begin{kid}
Throw ten paper balls at a target. The \textbf{mean} tells you where they landed on average (are you aiming too far left?). The \textbf{spread} (standard deviation) tells you how scattered they are. A great thrower has a mean right on the bullseye \emph{and} a small spread.
\end{kid}
